\documentclass[9pt,a4paper]{extarticle}
\usepackage[margin=1.3cm]{geometry}
\usepackage{fontspec}
\setmainfont{Times New Roman}
\usepackage{amsmath,amssymb}
\usepackage{booktabs,tabularx,array}
\usepackage{titlesec}
\usepackage{enumitem}

\titleformat{\section}{\normalsize\bfseries}{}{0em}{}[\vspace{-6pt}]
\titlespacing*{\section}{0pt}{4pt}{1pt}

\setlength{\parindent}{0pt}
\setlength{\parskip}{2pt}
\linespread{0.98}
\renewcommand{\arraystretch}{1.02}
\renewcommand{\tabularxcolumn}[1]{>{\raggedright\arraybackslash}p{#1}}
\pagestyle{empty}

\newcommand{\Cp}{\ensuremath{C_p}}
\newcommand{\CD}{\ensuremath{C_D}}
\newcommand{\CL}{\ensuremath{C_L}}
\newcommand{\ReD}{\ensuremath{\mathit{Re}}}
\newcommand{\blank}{\rule{1.5cm}{0.4pt}}

\begin{document}

{\large\bfseries Circular Cylinder --- Pre-lab Record}\hfill
{\small Engineering Practice 2 -- Aerodynamics 2 \quad RODH 115 (ex L50)}\\[1pt]
{\small Date \blank \qquad Group \blank \qquad Names \blank}

\section*{1 \ Objective}
Measure $\Cp(\theta)$ on 28 surface tappings at two speeds, with and without a trip wire;
integrate to get uncorrected $\CD'$, correct for blockage, and compare the integrated
$\CD{}_{\text{int}}$ against Roshko eq.~(2) $\CD{}_{\text{eq2}}$. Second aim: show that at the
same $\ReD$ the boundary-layer state alone moves separation from $\approx85^\circ$ to
$\approx130^\circ$ and $\CD$ from $\approx1.2$ to $\approx0.5$.

\section*{2 \ Apparatus (numbers to have at hand)}
$d = 102$\,mm, $b = 460$\,mm, $r = d/b = 0.2217$ (22\,\% blockage), 28 tappings + Pitot on
tubes 35/36, ethanol $\rho_e = 789$\,kg\,m$^{-3}$, manometer scale \blank\,mm/div,
engine settings 700 rpm ($U'\approx15.5$\,m/s, $\ReD'\approx1.0\times10^{5}$) and
1350 rpm ($U'\approx30.5$\,m/s, $\ReD'\approx2.1\times10^{5}$).
Hole 2 = $0^\circ$ (stagnation), holes 1 and 7 = $\pm30^\circ$ (mirror pair).
\par\vspace{3pt}
\noindent\fbox{\begin{minipage}[c][2.2cm][c]{\dimexpr\linewidth-2\fboxsep-2\fboxrule\relax}
\centering\small Sketch: tapping ring --- mark $0^\circ$, $\pm30^\circ$, $90^\circ$, $180^\circ$, and the wire position
\end{minipage}}

\section*{3 \ Procedure outline}
\begin{enumerate}[leftmargin=1.3em,itemsep=0pt,topsep=1pt]
\item Zero the manometer; check no bubbles; all tubes connected and numbered.
\item Read $T$ three times and $p_{\text{baro}}$ three times; note scale \blank\,mm/div.
\item Identify Pitot: the tube whose column matches hole 2 is \emph{total}; check
$q=\rho_e g B$ against $\tfrac12\rho U^2$ ($\approx146$\,Pa at 700, $\approx560$\,Pa at 1350).
\item Alignment: $\Cp(+30^\circ)$ vs $\Cp(-30^\circ)$. Benchmark $|\Delta\Cp{}_{30}|<0.01$
($\approx0.1^\circ$; the code only warns above $1.0^\circ$); if $0.01$--$0.12$, record it and
carry it into the $\theta$ error bar rather than repeating.
\item Run 1: 700 rpm, no wire. Settle; then \textbf{read all 30 columns in one pass,
three readings per tapping}.
\item Run 2: 1350 rpm, no wire. Same reading discipline.
\item Fit wire at \blank\,deg from stagnation (target 60--75$^\circ$); note its diameter.
\item Run 3: 700 rpm, tripped. \quad \item Run 4: 1350 rpm, tripped.
\item Optional run 5 if a third speed is offered.
\item Note anomalies per run (bubble, leak, drifting zero); keep raw heights only.
\end{enumerate}

\section*{4 \ Reduction, in this order}
\begin{enumerate}[leftmargin=1.3em,itemsep=0pt,topsep=1pt]
\item $\Cp = (h_i-h_s)/(h_t-h_s)$ --- $\rho$, $g$ cancel; it is a pure height ratio,
so $u(\Cp)=u(h)/B$.
\item $\CD' = \int_0^{\pi}\Cp\cos\theta\,\mathrm{d}\theta$;
$\ensuremath{C_{L,\text{asym}}}\approx \Delta\Cp{}_{30}/2$. Trapezoid on the actual tap spacing.
\item $V/V' = 1 + \tfrac14\CD' r + 0.82 r^2$ \quad (needs $\CD'$ first).
\item $\Cp^{\text{corr}} = 1 + (\Cp-1)/(V/V')^2$, then integrate again $\Rightarrow \CD{}_{\text{int}}$.
\item $\CD{}_{\text{eq2}} = (1-\tfrac12\CD' r - 2.5 r^2)\CD'$ (direct, eq.~2).
\item $\ReD' = \rho U' d/\mu$, $U_{\text{corr}} = U'(V/V')$; separation = angle where the
recovery ends and the plateau starts.
\end{enumerate}
$\CD{}_{\text{int}}$ and $\CD{}_{\text{eq2}}$ differ by 7--9\,\%: expected, not an error to hide.

\section*{5 \ Uncertainty plan}
Dominant term is one column reading divided by $B$: $u(\Cp)\approx0.027$ at the low speed
(1.9\,\%) and $0.007$ at the high speed (0.5\,\%) --- hence the high-speed runs are preferred
for quoted numbers. Tap angle tolerance $0.5^\circ/\sqrt3 = 0.29^\circ$; ring rotation from
the $\pm30^\circ$ pair (coherent over all taps). Separation angle carries $2.5^\circ$ when
the knee is sharp and $\pm20^\circ$ when the 12$^\circ$ slope window bottoms out on noise;
say which regime the run is in. Temperature and pressure enter only $\rho$ and $\mu$
(type B); they do not affect $\Cp$, which is a height ratio.

\section*{6 \ What I will record}
One row per tapping: \texttt{run, tripped, rpm, scale\_mm, Tair\_C\_1..3, pBaro\_Pa\_1..3,
hole, r\_1, r\_2, r\_3}. Then per run, for the logbook:
$\CD'$, $\CD{}_{\text{eq2}}$, $\CD{}_{\text{int}}$, $V/V'$, $\theta_{\text{sep}}$, $\ReD'$,
$u(\CD')$, $u(\theta_{\text{sep}})$, $\Delta\theta_{\text{ring}}$.

\begin{tabularx}{\textwidth}{@{}lXXXX@{}}
\toprule
Run & Wire & rpm & Purpose & Expected $\CD'$ \\
\midrule
1 & none & 700 & laminar baseline & $\approx1.1$ \\
2 & none & 1350 & same state, 4$\times$ $q$ & $\approx0.93$ \\
3 & fitted & 700 & tripped, low speed & $\approx0.46$ \\
4 & fitted & 1350 & tripped, high speed & $\approx0.60$ \\
5 & fitted & 1000 & mid-$Re$ tripped point & $\approx0.53$ \\
\bottomrule
\end{tabularx}

\section*{7 \ Questions I expect (and where I will check them)}
1 Vary $\ReD$ how, why, which speed more reliable (higher: $q\propto U^2$).
2 Best wire position (upstream of separation, 60--75$^\circ$).
3 Why separation differs at equal $\ReD$ (boundary-layer state).
4 Effect on drag (base pressure up, form drag down).
5 Equivalent $\ReD$ of the tripped runs ($\approx10^{6}$).
6 Two cases where tripping helps (golf ball; blade trip strip at low $\ReD$).
7 Blockage effect on $\Cp$ and $\CD$ (over-strong suction, high $\CD$).
8 Other error sources (reading, leaks, drift, shedding, 3D ends, tap count).
9 \textbf{Which Pitot is static --- verify at the bench by matching hole 2.}
10 Speed for a natural turbulent layer on a smooth cylinder ($\approx30$--$50$\,m/s, beyond
this tunnel). 11 Other surface features (roughness, dimples, trip dots, seams).

\end{document}
